\scanpage{192}
cas seulement, on peut choisir comme générateur de $B$ sur $A$
\[
 \Delta=\tau-\tau'
\]
soumis à la relation
\[
 \Delta^2=\Sigma^2-4\Pi\quad\in A
\]

On aura
\[
 \begin{aligned}
 \tau&=\tfrac12(\Delta+\Sigma)\\
 \tau'&=\tfrac12(-\Delta+\Sigma)
 \end{aligned}
\]

Mais en car. 2, on aura
\[
 \Delta=\tau+\tau'=\Sigma\in A\qquad\text{(en car. 2)}
\]
donc $\Delta$ ne risque pas d'engendrer $B$ sur $A$ en car. 2, donc pas non plus si $2$ non inv. dans $k$ i.e. $k/2k\ne0$.

\begin{center}\rule{2.4cm}{0.3pt}\end{center}

Cas de quatre variables $\unsure{x_1,y_1,z_1,t_1|x,y,z,t}$.

Soit $B_4(\mathcal A)\subset\mathrm{Sym}(\overset{\vee}{\mathcal A}^{4})$ l'anneau des fonctions polynômes $F(x,y,z,t)$
% archive sheet number: 191
